\documentclass[a4paper,12pt]{article}
\usepackage{amssymb}

\title{AlDa-Klausur 1}

\begin{document}
\maketitle
\section*{Aufgabe 1: O-Kalkül}
\subsection*{Aufgabe 1.a: O-Notation}
\subsubsection*{(1)}
$$O(\frac{log n}{log log n}) \subset O((log n)^2) \subset O(\sqrt{n^{1.5}} + 1000 \cdot n^{1.2}) \subset O(2^{\sqrt{n}})$$
\subsubsection*{(2)}
\subsection*{Aufgabe 1.b: Laufzeitabschätzung}
\subsubsection*{(1)}
int $$k = 0 \rightarrow O(1)$$
outer loop: for i in 1..n, double i $$\rightarrow \frac{n}{2} \rightarrow O(log n)$$
inner loop: for j in n..i, sub j 1 $$\rightarrow i := n, j := i \rightarrow O(n)$$
Gesamt: outer loop * inner loop: $$O(n log n)$$

\subsubsection*{(2)}
outer loop: for i in 0..n, add i 2 => n/2 Steps => n => O(n)
inner loop: for j in n..1, sub j 1 => n => O(n)
Gesamt: $$O(n^2)$$
\subsubsection*{(3)}
Queue FIFO: n
for v in a => O(n)
while not queue <- empty()? => worst Case O(n)
while-nested for loop: for w : v <- succ => if := O(1), for := O(n)
Gesamt: $$O(n) + O(n^2) \rightarrow O(n^2)$$
\section*{Aufgabe 2: Teile und Herrsche}
\subsection*{Aufgabe 2.a: Laufzeit von Teile und Herrsche}
\subsubsection*{(1)}
int d ... => O(1)
if d == (0 or 1)? => O(1)
int m ... => O(1)
for i in left..m, add i 1 => O(m) => O(n/2)
return ... => Watershed-Konstante mit Master-Theorem => Watershed := 1 => $$n^1$$ => O(n)
Gesamt: O(3) => O(1) => O(1) + O(n/2) + O(log n) => O(n log n)
\subsubsection*{(2)}
if ... => O(1)
m $$\rightarrow O(1)$$
h(...);
h(...);
=> Watershed und Master-Theorem $$\rightarrow \Theta(n)$$
\subsection*{Aufgabe 2.b: Binärdarstellung}
\subsubsection*{(1)}
$$binary(x\in \mathbb{N}) := $$\\
n $<$ 2 ? ret n\\
else ret binary(div n 2) ++ print mod n 2
\subsubsection*{(2)}
Master-Theorem: O(log n)
\subsection*{Aufgabe 2.c: Anzahl von Vertauschungen}
\subsubsection*{(1)}
$dac\_mergesort(a :: [n \in \mathbb{Z}], l, m, r \in \mathbb{Z}) :=$\\
\indent $i, j, k \in \mathbb{Z}, aux :: [x \in \mathbb{Z}]$\\ 
\indent for i in m+1..l, sub i 1 $\rightarrow$ aux (i-1) := a (i-1)\\
\indent for j in m..r, add j 1 $\rightarrow$ aux (r+m-j) := a (j+1)\\
\indent n $\leftarrow$ 0\\
\indent\indent for k in l..r, add k 1\\
\indent\indent\indent if aux (j) $<$ aux (i) $\rightarrow$ a (k) := aux (j); sub j 1; \\
\indent\indent\indent\indent if i $\le$ m $\rightarrow$ n $\leftarrow$ add (n (m-i+1))\\
\indent\indent\indent else $\rightarrow$ a (k) := aux (i); add i 1\\
\indent ret n
\\
\\$dac\_count(a :: [n \in \mathbb{Z}], l, r \in \mathbb{Z}) :=$\\
\indent if r $\le$ l $\rightarrow$ ret Nil\\
\indent else\\
\indent\indent m $\leftarrow$ div (add (l r) 2)\\
\indent\indent dac\_count(a, l, m)\\
\indent\indent dac\_count(a, m+1, r)\\
\indent\indent dac\_mergesort(a, l, m, r)

\subsubsection*{(2)}
mergesort first loop: Copy-Aufruf: O(n)\\
mergesort second loop: Copy-Aufruf: O(n)\\
$\rightarrow$ O(n)\\
third loop: O(n)\\
mergesort :: O(n)\\
count :: Master-Theorem :: O(log n)\\
Gesamt: $O(n log n)$

\section*{Aufgabe 3: Hashen}
\subsection*{Aufgabe 3.a: Multiplikative Methode}

\section*{Aufgabe 4: Sortieren}
\subsection*{Aufgabe 4.a: Mergesort}
\paragraph{50, 51, 2, 9, -2, 17}



\end{document}